\chapter{CONCEPTION DE LA SOLUTION}
\label{ch:conception}

Ce chapitre présente l'architecture cible de la Cloud Data Platform. La conception reprend la logique en sept couches décrite dans les spécifications, puis la projette sur AWS EKS, Kubernetes et les artefacts du dépôt.

%------------------------------------------------------------------------------
\section{Architecture globale}
%------------------------------------------------------------------------------

La plateforme comporte sept couches logiques : sources, ingestion, stockage, calcul, virtualisation, présentation et observabilité. Elle est hébergée sur \textbf{AWS EKS}~\cite{aws-eks} en environnement \texttt{dev}, avec un design compatible avec une évolution multi-cloud ou on-premise grâce à Kubernetes~\cite{kubernetes}, MinIO~\cite{minio} et Iceberg~\cite{apache-iceberg}.

La figure~\ref{fig:arch-globale} représente ces couches et met en évidence le chemin suivi par les données jusqu'aux usages analytiques.

\begin{figure}[H]
\centering
\maybeimage[width=0.95\textwidth,height=0.72\textheight,keepaspectratio]{images/screenshots_pipeline/architecture_globale.png}
\caption{Architecture logique en couches de la Cloud Data Platform}
\label{fig:arch-globale}
\end{figure}

Cette représentation donne une lecture fonctionnelle de la plateforme. Les données entrent par la couche sources, sont captées par les services d'ingestion, puis déposées dans le stockage objet où elles sont structurées selon les zones Bronze, Silver et Gold. Les traitements de calcul assurent ensuite les transformations et la préparation des tables analytiques, tandis que la couche de virtualisation expose ces jeux de données aux outils de restitution. L'intérêt de cette vue est de montrer que chaque composant technique occupe une place précise dans le parcours de la donnée : Kafka et NiFi se concentrent sur l'acquisition, MinIO et Iceberg sur la persistance, Spark et dbt sur la transformation, Dremio et Power BI sur la consommation. La couche d'observabilité encadre l'ensemble afin de suivre l'état des services et de détecter les incidents.

Les captures détaillées d'exécution sont présentées dans le chapitre de réalisation. Dans ce chapitre, seules les illustrations qui aident à comprendre les choix d'architecture sont conservées.

%------------------------------------------------------------------------------
\section{Déploiement physique sur AWS EKS}
%------------------------------------------------------------------------------

\subsection{Landing zone}

La landing zone est composée d'un VPC multi-AZ, de sous-réseaux privés et publics, d'un NAT Gateway unique en environnement \texttt{dev}, de clés KMS dédiées et d'un cluster EKS managé. Les modules Terraform \texttt{vpc}, \texttt{kms}, \texttt{iam}, \texttt{eks}, \texttt{ecr} et \texttt{lambda-scaling} isolent les responsabilités et permettent une évolution contrôlée.

La figure~\ref{fig:conception-socle-aws-eks} illustre les principales briques AWS de cette landing zone : réseau, routage et cluster EKS.

\begin{figure}[H]
\centering
\PFEScreenshotItem{Vue AWS du VPC dédié à la plateforme}{images/screenshots_infra/aws_vpc_overview.png}
\PFEScreenshotItem{Organisation des sous-réseaux et des tables de routage}{images/screenshots_infra/aws_subnets_routes.png}
\PFEScreenshotItem{Vue générale du cluster EKS \texttt{dataplatform-dev}}{images/screenshots_infra/aws_eks_cluster_overview.png}
\caption{Socle cible AWS/EKS retenu pour la conception}
\label{fig:conception-socle-aws-eks}
\end{figure}

La figure regroupe trois vues complémentaires du socle physique. La première montre le VPC dédié, qui constitue l'enveloppe réseau de la plateforme et permet d'isoler les ressources du reste du compte AWS. La deuxième détaille les sous-réseaux et les tables de routage : les sous-réseaux publics portent les composants d'accès contrôlé, alors que les sous-réseaux privés hébergent les nœuds Kubernetes et les services applicatifs. La troisième vue présente le cluster EKS \texttt{dataplatform-dev}, point d'exécution principal des workloads. Cette lecture met en évidence le choix d'une infrastructure managée : AWS fournit les fondations réseau et le plan de contrôle Kubernetes, tandis que la plateforme garde la maîtrise de ses composants data via Terraform et les manifests Kubernetes.

\subsection{Node groups}

Deux groupes de nœuds séparent les contraintes stateful et compute :
\begin{itemize}
    \item \texttt{stateful-ng} : workloads persistants comme MinIO, PostgreSQL, Hive et Kafka ;
    \item \texttt{compute-ng} : workloads élastiques comme Spark, Airflow, Dremio, observabilité et jobs ponctuels.
\end{itemize}

Cette séparation rend possible le \textbf{scale-to-zero} du compute sans arrêter les services de stockage, ce qui répond à la contrainte d'optimisation des coûts.

%------------------------------------------------------------------------------
\section{Stratégie des environnements DEV, UAT et PROD}
%------------------------------------------------------------------------------

La plateforme est conçue selon une progression d'environnements qui sépare les usages de développement, de validation métier et d'exploitation. Cette séparation évite de mélanger les expérimentations techniques avec les données et services destinés aux utilisateurs finaux.

La figure~\ref{fig:environments-dev-uat-prod} présente cette progression. Le tableau~\ref{tab:environnements} précise le rôle et le niveau d'exigence associé à chaque environnement.

\begin{figure}[H]
\centering
\maybeimage[width=0.96\textwidth,height=0.58\textheight,keepaspectratio]{images/screenshots_pipeline/environments.png}
\caption{Organisation cible des environnements DEV, UAT et PROD}
\label{fig:environments-dev-uat-prod}
\end{figure}

Cette figure illustre le cycle de promotion prévu pour la plateforme. L'environnement \textbf{DEV} sert d'espace de construction : les modules d'infrastructure, les déploiements Kubernetes et les traitements de données y sont modifiés rapidement avec un niveau de risque limité. L'environnement \textbf{UAT} joue le rôle de filtre de validation ; il permet de vérifier que les flux fonctionnent de bout en bout et que les résultats correspondent aux attentes avant toute mise en production. L'environnement \textbf{PROD} représente enfin la cible d'exploitation, avec des paramètres plus stricts de disponibilité, de sécurité et de supervision. La séparation visible sur la figure évite qu'un changement expérimental influence directement les usages finaux.

\begin{table}[H]
\centering
\caption{Rôle des environnements de la plateforme}
\label{tab:environnements}
\small
\begin{tabular}{@{}p{0.16\textwidth}p{0.68\textwidth}@{}}
\toprule
\textbf{Environnement} & \textbf{Rôle dans la conception} \\
\midrule
\textbf{DEV} & Environnement de construction et d'expérimentation. Il sert à développer les modules Terraform, les manifests Kubernetes, les DAGs Airflow, les jobs Spark et les modèles dbt avec des ressources réduites et des données de test. \\
\textbf{UAT} & Environnement de recette. Il reproduit la chaîne fonctionnelle de bout en bout pour valider les flux, les tables Silver/Gold, les contrôles de qualité, les accès aux interfaces et les scénarios de démonstration avant promotion. \\
\textbf{PROD} & Environnement d'exploitation. Il applique des paramètres plus stricts : haute disponibilité, quotas adaptés, sauvegardes, supervision renforcée, secrets réels, contrôle d'accès plus fin et procédures de changement maîtrisées. \\
\bottomrule
\end{tabular}
\end{table}

Dans cette logique, un changement est d'abord testé en \textbf{DEV}, stabilisé en \textbf{UAT}, puis promu vers \textbf{PROD}. Les mêmes principes d'infrastructure as code, de namespaces, de secrets et de CI/CD sont conservés entre les environnements, mais les tailles de ressources, les règles d'accès et les paramètres de résilience évoluent selon le niveau de criticité.

%------------------------------------------------------------------------------
\section{Stratégie des namespaces Kubernetes}
%------------------------------------------------------------------------------

Le tableau~\ref{tab:namespaces} décrit la segmentation Kubernetes retenue. Chaque namespace regroupe des workloads ayant une responsabilité technique homogène.

\begin{table}[H]
\centering
\caption{Namespaces \texttt{platform-*} alignés avec le cahier des charges}
\label{tab:namespaces}
\small
\begin{tabular}{@{}p{0.26\textwidth}p{0.6\textwidth}@{}}
\toprule
\textbf{Namespace} & \textbf{Workloads prévus} \\
\midrule
\texttt{platform-ingestion} & Strimzi Kafka, Kafka Connect, Debezium, NiFi \\
\texttt{platform-storage} & MinIO Operator, MinIO Tenant, PostgreSQL source \\
\texttt{platform-metastore} & Hive Metastore, PostgreSQL HMS \\
\texttt{platform-compute} & Spark Operator, SparkApplication, dbt jobs \\
\texttt{platform-orchestr} & Airflow, base metadata Airflow \\
\texttt{platform-serving} & Dremio, services SQL/BI \\
\texttt{platform-monitoring} & Prometheus, Alertmanager, Grafana \\
\texttt{platform-logging} & Fluent Bit, OpenObserve \\
\texttt{platform-security} & cert-manager~\cite{cert-manager}, Cloudflare Operator, secrets d'exposition \\
\bottomrule
\end{tabular}
\end{table}

Chaque namespace reçoit une \texttt{ResourceQuota}, une \texttt{NetworkPolicy} de base et un niveau Pod Security Admission adapté. Cette stratégie limite le blast radius et facilite l'explication de l'architecture au jury.

%------------------------------------------------------------------------------
\section{Conception du pipeline médaillon}
%------------------------------------------------------------------------------

Le pipeline est conçu autour d'un flux principal PostgreSQL $\rightarrow$ Kafka $\rightarrow$ Iceberg, complété par une ingestion fichiers. La séparation Bronze/Silver/Gold permet d'organiser la qualité et la finalité des données.

La figure~\ref{fig:flux-medallion} illustre le flux cible, tandis que le tableau~\ref{tab:medallion} explicite le rôle de chaque couche de données.

\begin{figure}[H]
\centering
\makebox[\textwidth][c]{%
  \maybeimage[width=1.08\textwidth,height=0.55\textheight,keepaspectratio]{images/pipeline_data.png}%
}
\caption{Flux de données médaillon cible}
\label{fig:flux-medallion}
\end{figure}

La figure met l'accent sur la transformation progressive de la donnée. Les sources opérationnelles, comme PostgreSQL, et les fichiers déposés alimentent d'abord la zone \textbf{Bronze}, qui conserve les données dans un état proche de l'origine afin de préserver la traçabilité. Les traitements Spark et dbt produisent ensuite la zone \textbf{Silver}, où les données sont nettoyées, typées, filtrées et préparées pour les jointures métier. La zone \textbf{Gold} contient les tables les plus proches des usages décisionnels : indicateurs consolidés, agrégats et vues prêtes pour la BI. Cette organisation réduit le couplage entre ingestion, transformation et restitution, car chaque couche possède un rôle clair et peut être contrôlée séparément.

\clearpage

\begin{table}[H]
\centering
\caption{Rôle des couches médaillon}
\label{tab:medallion}
\small
\begin{tabular}{@{}p{0.18\textwidth}p{0.68\textwidth}@{}}
\toprule
\textbf{Couche} & \textbf{Rôle} \\
\midrule
Bronze & Données brutes, historisées, proches de la source, écrites en Iceberg pour conserver la traçabilité. \\
Silver & Données nettoyées, typées, dédupliquées et prêtes pour les jointures métier. \\
Gold & Agrégats orientés analyse, indicateurs stabilisés et tables consommables par BI. \\
\bottomrule
\end{tabular}
\end{table}

La couche d'ingestion conserve deux modes complémentaires. NiFi permet de représenter les flux semi-structurés et les scénarios d'intégration métier, tandis qu'Airflow pilote les traitements batch et dbt. Les captures concrètes de ces interfaces sont placées dans le chapitre de réalisation, car elles servent de preuve d'implémentation.

Le stockage objet est organisé par zones Bronze, Silver et Gold. Cette organisation rend visible la progression des données depuis leur état brut jusqu'aux tables analytiques.

%------------------------------------------------------------------------------
\section{Sécurité et gouvernance technique}
%------------------------------------------------------------------------------

La sécurité est intégrée dès la conception :
\begin{itemize}
    \item \textbf{identité} : OIDC GitHub vers AWS et IRSA pour les controllers Kubernetes ;
    \item \textbf{chiffrement} : clés KMS distinctes pour EBS, secrets et logs ;
    \item \textbf{réseau} : API EKS privée par défaut, NetworkPolicies et exposition UI via Cloudflare Tunnel ;
    \item \textbf{secrets} : templates versionnés, valeurs réelles générées localement et exclues de Git ;
\end{itemize}

Les captures AWS et Cloudflare correspondantes sont déplacées dans la partie réalisation afin de montrer l'application concrète de ces principes.

%------------------------------------------------------------------------------
\section{Observabilité, logging et dashboarding}
%------------------------------------------------------------------------------

L'observabilité est conçue comme une couche transversale, et non comme un ajout final. Les métriques Kubernetes et applicatives sont collectées par Prometheus ; Grafana fournit les dashboards techniques ; Fluent Bit route les logs vers OpenObserve ; les ServiceMonitors ciblent les composants critiques.

Le tableau~\ref{tab:observabilite-conception} synthétise cette conception en distinguant métriques, logs, dashboards et alertes.

\begin{table}[H]
\centering
\caption{Conception de l'observabilité}
\label{tab:observabilite-conception}
\small
\begin{tabular}{@{}p{0.22\textwidth}p{0.64\textwidth}@{}}
\toprule
\textbf{Pilier} & \textbf{Mise en œuvre prévue} \\
\midrule
Métriques & kube-prometheus-stack, métriques Kubernetes, Spark Operator, MinIO, CNPG, Strimzi, Airflow et Dremio. \\
Logs & Fluent Bit en DaemonSet, enrichissement Kubernetes, sortie OpenObserve, recherche par namespace/pod/job. \\
Dashboards & Grafana pour santé cluster, ressources, Kafka, PostgreSQL, Spark, Airflow et stockage ; Power BI pour indicateurs métier Gold. \\
Alerting & Alertmanager pour alertes techniques : pods crashloop, PVC presque pleins, jobs échoués, connecteurs CDC indisponibles. \\
\bottomrule
\end{tabular}
\end{table}

%------------------------------------------------------------------------------
\section{Organisation technique des livrables}
%------------------------------------------------------------------------------

\begin{lstlisting}[language=bash,caption={Organisation technique des livrables}]
plateforme-data/
  terraform/modules/{vpc,kms,iam,eks,ecr,lambda-scaling}
  terraform/environments/dev/
  bootstrap/phase-2..phase-10/
  docker/spark-iceberg/
  dags/
  scripts/bootstrap-phase*.sh
  docs/specs/, docs/phases/, docs/02-architecture-decision-records.md
  rapport/
\end{lstlisting}

Les captures des dossiers, registres, buckets, UIs et validations terminal sont volontairement regroupées dans le chapitre suivant, car elles décrivent la mise en œuvre réelle.

\section*{Conclusion du chapitre}
\addcontentsline{toc}{section}{Conclusion du chapitre 4}

Ce chapitre a traduit les exigences en une architecture cible structurée. La solution repose sur sept couches logiques : sources, ingestion, stockage, calcul, virtualisation, présentation et observabilité. Cette décomposition permet d'isoler les responsabilités, de faire évoluer les composants indépendamment et de conserver une lecture claire du parcours de la donnée.

Le socle physique retenu s'appuie sur AWS EKS et sur une landing zone décrite par Terraform. Le VPC multi-AZ, les sous-réseaux privés et publics, le NAT Gateway, les clés KMS et les rôles IAM fournissent les fondations réseau et sécurité. La séparation entre node groups stateful et compute répond à un double objectif : préserver la disponibilité des services persistants et permettre une réduction des coûts lorsque les traitements élastiques ne sont pas utilisés.

La conception prend également en compte le cycle de vie de la plateforme. Les environnements DEV, UAT et PROD distinguent construction, recette et exploitation. Les namespaces Kubernetes regroupent les workloads par fonction : ingestion, stockage, métastore, calcul, orchestration, serving, monitoring, logging et sécurité. Cette organisation facilite l'application de quotas, de politiques réseau et de règles d'accès adaptées à chaque périmètre.

Le pipeline médaillon formalise ensuite le chemin des données. PostgreSQL, Kafka et les flux fichiers alimentent Bronze ; Spark et dbt assurent les transformations Silver et Gold ; Dremio expose les tables analytiques à la BI. En parallèle, Prometheus, Grafana, Fluent Bit et OpenObserve constituent la couche transverse d'observabilité. La sécurité est intégrée dès la conception au travers d'OIDC, IRSA, KMS, secrets hors Git et exposition contrôlée par Cloudflare Tunnel.

Enfin, l'organisation technique des livrables montre que cette conception est directement exploitable : modules Terraform, manifests Kubernetes, scripts de bootstrap, image Spark, DAGs et documentation sont structurés pour accompagner le déploiement par phases.

Le chapitre suivant confronte cette architecture à sa mise en œuvre réelle. Il présente les composants déployés, les preuves d'exécution, les validations réalisées et les éléments restant à finaliser pour renforcer l'industrialisation.

\clearpage
